% !TeX root = main.tex

\chapter{Wasserstein Gradient Flows}
\section{Motivation}

In Euclidean optimization, gradient descent evolves points according to:

\[
\dot x_t = -\grad f(x_t).
\]

The fundamental idea of Wasserstein gradient flows is the following:

\begin{quote}
Instead of evolving points in Euclidean space, evolve probability distributions in Wasserstein space.
\end{quote}

The resulting framework connects:

\begin{itemize}
    \item partial differential equations,
    \item diffusion processes,
    \item entropy dissipation,
    \item stochastic control,
    \item generative modeling,
    \item reinforcement learning.
\end{itemize}

This chapter develops the formalism of Wasserstein gradient flows (WGF), emphasizing:

\begin{enumerate}[label=(\roman*)]
    \item variational principles,
    \item continuity equations,
    \item free energy functionals,
    \item the Jordan--Kinderlehrer--Otto (JKO) scheme,
    \item connections to diffusion and learning.
\end{enumerate}

\section{Gradient Flows in Euclidean Space}

We briefly recall ordinary gradient descent.

\subsection{Gradient descent ODE}

Let
\[
f : \R^d \to \R
\]
be smooth.

The gradient flow equation is

\[
\dot x_t = -\grad f(x_t).
\]

This equation follows the direction of steepest descent.

\subsection{Energy dissipation}

Along solutions:

\[
\frac{d}{dt}f(x_t)
=
\grad f(x_t) \cdot \dot x_t.
\]

Substituting the gradient flow equation gives:

\[
\frac{d}{dt}f(x_t)
=
-\|\grad f(x_t)\|^2
\le 0.
\]

Thus gradient flows dissipate energy monotonically.

\subsection{Metric interpretation}

Gradient descent depends on the geometry of the underlying space.

In Euclidean space, the metric tensor is the identity.

In Wasserstein geometry, the metric structure is different, leading to different gradient flows.

\section{Functionals on Probability Space}

Wasserstein gradient flows evolve probability measures to decrease an energy functional.

\subsection{Energy functionals}

A functional is a map:

\[
\mathcal F : \cP_2(\R^d) \to \R.
\]

Important examples include:

\begin{enumerate}[label=(\alph*)]
    \item potential energies,
    \item entropy functionals,
    \item interaction energies,
    \item KL divergences.
\end{enumerate}

\subsection{Potential energy}

Given a confining potential
\[
V : \R^d \to \R,
\]
define:

\[
\mathcal V(\rho)
=
\int V(x)d\rho(x).
\]

\subsection{Entropy}

For a density $\rho$:

\[
\mathcal H(\rho)
=
\int \rho(x)\log \rho(x)dx.
\]

Note the sign convention differs from Shannon entropy.

\subsection{Free energy}

A central object is the free energy:

\[
\mathcal F(\rho)
=
\int V(x)d\rho(x)
+
\beta^{-1}
\int \rho(x)\log \rho(x)dx.
\]

This combines:

\begin{itemize}
    \item energy minimization,
    \item entropy regularization.
\end{itemize}

The parameter:

\[
\beta^{-1}
\]
plays the role of temperature.

\section{First Variations}

To define gradient flows on probability space we require a notion of derivative.

\subsection{Functional derivatives}

\begin{definition}[First variation]
Let $\mathcal F$ be a functional defined on a class of probability densities.
The first variation of $\mathcal F$ at $\rho$ is a function
\[
\frac{\delta \mathcal F}{\delta \rho}(\rho)
\]
such that, for every smooth perturbation $\eta$ satisfying
\[
\int_{\R^d}\eta(x)\,dx=0,
\]
one has
\[
\left.
\frac{d}{d\varepsilon}
\mathcal F(\rho+\varepsilon\eta)
\right|_{\varepsilon=0}
=
\int_{\R^d}
\frac{\delta \mathcal F}{\delta \rho}(\rho)(x)\,\eta(x)\,dx.
\]
\end{definition}

The zero-mass condition on $\eta$ ensures that $\rho+\varepsilon\eta$ remains a probability density to first order.
On probability densities, the first variation is defined only up to an additive constant, since constants vanish against zero-mass perturbations.

\subsection{Examples}

\paragraph{Potential energy}

For
\[
\mathcal V(\rho)=\int V d\rho,
\]
one has:

\[
\frac{\delta \mathcal V}{\delta \rho}=V.
\]

\paragraph{Entropy}

For
\[
\mathcal H(\rho)
=
\int \rho\log \rho,
\]
one obtains:

\[
\frac{\delta \mathcal H}{\delta \rho}
=
\log \rho +1.
\]

\paragraph{Free energy}

For
\[
\mathcal F(\rho)
=
\int V d\rho
+
\beta^{-1}
\int \rho\log \rho,
\]
we get:

\[
\frac{\delta \mathcal F}{\delta \rho}
=
V + \beta^{-1}(\log \rho +1).
\]

\section{Wasserstein Gradient Flows}

\subsection{Formal definition}

The Wasserstein gradient flow of a functional $\mathcal F$ is formally:

\[
\explanationlink{wasserstein_gradient_flow_derivation}{\partial_t \rho_t
=
\divergence
\left(
\rho_t
\grad
\frac{\delta \mathcal F}{\delta \rho}(\rho_t)
\right)}.
\]

Equivalently:

\[
\partial_t \rho_t
+
\divergence(\rho_t v_t)=0,
\]

with velocity field:

\[
v_t
=
-\grad
\frac{\delta \mathcal F}{\delta \rho}(\rho_t).
\]

\subsection{Interpretation}

Compare with Euclidean gradient flow:

\[
\dot x_t = -\grad f(x_t).
\]

In Wasserstein space:

\begin{itemize}
    \item points become distributions,
    \item velocities become transport vector fields,
    \item gradients become functional derivatives.
\end{itemize}

\section{The Fokker--Planck Equation}

One of the most important examples of a Wasserstein gradient flow is the Fokker--Planck equation.

\subsection{Langevin dynamics}

Consider the stochastic differential equation:

\[
dX_t
=
-\grad V(X_t)dt
+
\sqrt{2\beta^{-1}}dW_t.
\]

Its probability density evolves according to:

\[
\partial_t \rho_t
=
\divergence(\rho_t \grad V)
+
\beta^{-1}\Delta \rho_t.
\]

\subsection{Gradient flow structure}

\begin{theorem}
The Fokker--Planck equation is the Wasserstein gradient flow of the free energy:

\[
\mathcal F(\rho)
=
\int Vd\rho
+
\beta^{-1}
\int \rho\log \rho.
\]
\end{theorem}

\subsection{Derivation}

Using:

\[
\frac{\delta \mathcal F}{\delta \rho}
=
V+\beta^{-1}(\log \rho+1),
\]
we compute:

\begin{align*}
\partial_t \rho
&=
\divergence
\left(
\rho
\grad
(V+\beta^{-1}\log \rho)
\right)
\\
&=
\divergence(\rho\grad V)
+
\beta^{-1}
\divergence(\rho\grad \log \rho).
\end{align*}

Since
\[
\rho\grad \log \rho = \grad \rho,
\]
we obtain:

\[
\partial_t \rho
=
\divergence(\rho\grad V)
+
\beta^{-1}\Delta \rho.
\]

\section{Energy Dissipation}

Gradient flows dissipate free energy.

\begin{theorem}[Energy dissipation]
Along Wasserstein gradient flows:

\[
\frac{d}{dt}\mathcal F(\rho_t)
=
-
\int
\left\|
\grad
\frac{\delta \mathcal F}{\delta \rho}
\right\|^2
 d\rho_t.
\]
\end{theorem}

Thus:

\[
\frac{d}{dt}\mathcal F(\rho_t) \le 0.
\]

\begin{remark}
This is the infinite-dimensional analogue of:

\[
\frac{d}{dt}f(x_t)
=
-\|\grad f(x_t)\|^2.
\]
\end{remark}

\section{The Jordan--Kinderlehrer--Otto Scheme}

The JKO scheme gives a variational discretization of Wasserstein gradient flows.

\subsection{Implicit Euler in Wasserstein space}

Recall ordinary gradient descent:

\[
x_{k+1}
=
\argmin_x
\left(
\frac{1}{2\tau}\|x-x_k\|^2
+
f(x)
\right).
\]

The JKO scheme replaces Euclidean distance with Wasserstein distance.

\begin{definition}[JKO scheme]
Given timestep $\tau>0$, define recursively:

\[
\rho_{k+1}
=
\argmin_\rho
\left(
\frac{1}{2\tau}
W_2^2(\rho,\rho_k)
+
\mathcal F(\rho)
\right).
\]
\end{definition}

\subsection{Interpretation}

Each step balances:

\begin{itemize}
    \item staying close to the previous distribution,
    \item decreasing free energy.
\end{itemize}

Thus Wasserstein gradient flows emerge from repeated proximal optimization.

\subsection{Conceptual importance}

The JKO scheme is foundational because it:

\begin{itemize}
    \item defines gradient flows variationally,
    \item avoids explicit PDE analysis,
    \item links optimization and transport geometry.
\end{itemize}

Many modern ML methods can be interpreted as approximate JKO steps.

\section{KL Divergence as Free Energy}

Suppose:

\[
\pi(x)
\propto
e^{-V(x)}.
\]

Then:

\[
\KL(\rho\|\pi)
=
\int \rho\log \rho
+
\int Vd\rho
+
\text{constant}.
\]

Thus:

\begin{quote}
KL divergence is a free energy.
\end{quote}

Consequently:

\begin{quote}
The Fokker--Planck equation is gradient descent on KL divergence in Wasserstein space.
\end{quote}

This observation is central for:

\begin{itemize}
    \item diffusion models,
    \item score-based modeling,
    \item entropy-regularized RL,
    \item Schr\"odinger bridges.
\end{itemize}

\section{Connections to Generative Modeling}

\subsection{Diffusion models}

Diffusion models gradually transform noise distributions into data distributions.

The forward diffusion process satisfies a Fokker--Planck equation.

The reverse process can be interpreted as transport in probability space.

\subsection{Flow models}

Continuous normalizing flows evolve densities according to:

\[
\partial_t \rho_t
+
\divergence(\rho_t v_t)=0.
\]

Thus they directly implement deterministic transport.

\subsection{Drifting models}

Drifting models train generators whose induced distributions evolve during optimization.

A key question is:

\begin{quote}
Does SGD in parameter space induce approximate Wasserstein gradient flow in distribution space?
\end{quote}

This perspective motivates drifting generative modeling.

\section{Connections to Reinforcement Learning}

Policies are conditional distributions:

\[
\pi(a|s).
\]

Policy optimization may be interpreted as transport in action-distribution space.

Instead of:

\[
\theta_t \to \theta_{t+1},
\]

one may study:

\[
\pi_t \to \pi_{t+1}.
\]

Recent work formulates RL as Wasserstein gradient flow toward soft optimal policies.

\subsection{Entropy-regularized RL}

Soft RL objectives often involve:

\[
\mathbb E[-Q(a,s)]
+
\beta^{-1}\KL(\pi\|\pi_0).
\]

This resembles free-energy minimization.

Consequently:

\begin{quote}
Soft policy optimization may be viewed as gradient flow in probability space.
\end{quote}

This viewpoint underlies drifting-field policies.

\section{Otto Calculus Revisited}

Recall the formal metric:

\[
\langle \dot \rho_1,\dot \rho_2 \rangle_\rho
=
\int \grad \phi_1 \cdot \grad \phi_2 d\rho.
\]

Under this geometry:

\[
\partial_t \rho
=
-\mathrm{grad}_{W_2}\mathcal F(\rho).
\]

Thus Wasserstein gradient flow is literally gradient descent under the Wasserstein metric.

This is the conceptual core of the theory.

\section{Summary}

The main ideas of this phase are:

\begin{enumerate}[label=(\arabic*)]
    \item Wasserstein space admits gradient-flow dynamics.
    \item Functionals on probability space define evolution equations.
    \item The Fokker--Planck equation is a Wasserstein gradient flow.
    \item Free-energy minimization combines energy and entropy.
    \item The JKO scheme defines variational time discretization.
    \item KL divergence can be interpreted as a free energy.
    \item Many modern AI models evolve distributions rather than points.
\end{enumerate}

The key conceptual shift is:

\begin{quote}
Optimization algorithms may be interpreted as transport dynamics in probability space.
\end{quote}

This idea will be fundamental for:

\begin{itemize}
    \item Schr\"odinger bridges,
    \item entropy-regularized optimal transport,
    \item diffusion models,
    \item score-based generative modeling,
    \item Wasserstein reinforcement learning,
    \item drifting-field policies.
\end{itemize}

In the next phase we study policy optimization and reinforcement learning directly in Wasserstein space.

\section*{Recommended Reading}

\begin{enumerate}
    \item L. Ambrosio, N. Gigli, and G. Savar\'e, \emph{Gradient Flows in Metric Spaces and in the Space of Probability Measures}.
    \item F. Otto, ``The Geometry of Dissipative Evolution Equations.''
    \item R. Jordan, D. Kinderlehrer, and F. Otto, ``The Variational Formulation of the Fokker--Planck Equation.''
    \item C. Villani, \emph{Optimal Transport: Old and New}.
    \item F. Santambrogio, \emph{Optimal Transport for Applied Mathematicians}.
\end{enumerate}
